% ------------------------------------------------------------------------
% ------------------------------------------------------------------------
% ICMC: Modelo de Trabalho Acadêmico (tese de doutorado, dissertação de
% mestrado e trabalhos monográficos em geral) em conformidade com 
% ABNT NBR 14724:2011: Informação e documentação - Trabalhos acadêmicos -
% Apresentação
% ------------------------------------------------------------------------
% ------------------------------------------------------------------------

% Opções: 
%   Qualificação         = qualificacao 
%   Curso                = doutorado/mestrado
%   Situação do trabalho = pre-defesa/pos-defesa (exceto para qualificação)
% -- opções do pacote babel --
% Idioma padrão = brazil
	%spanish,			% idioma adicional para hifenização
	%english,			% idioma adicional para hifenização
	%brazil				% o último idioma é o principal do documento
\documentclass[mestrado, spanish, english, brazil]{packages/icmc}

% ---
% Pacotes Opcionais
% ---
\usepackage{rotating}           % Usado para rotacionar o texto
\usepackage[all,knot,arc,import,poly]{xy}   % Pacote para desenhos gráficos
% Este pacote pode conflitar com outros pacotes gráficos como o ``pictex''
% Então é necessário usar apenas um dos pacotes conflitantes

\newcommand{\VerbL}{0.52\textwidth}
\newcommand{\LatL}{0.42\textwidth}


% ---
% Informações de dados para CAPA e FOLHA DE ROSTO
% ---
\titulo{Título do TCC}
\autor[Neuhaus, L. A.]{Lucas Arendt Neuhaus}
\orientador[Orientador]{Prof. Dr.}{Jefferson Lins da Silva}
%\coorientador{Prof. Dr.}{Fulano de Tal}
\curso{CCMC}
\data{5}{9}{2018} % Data do depósito
% ---


% ---
% RESUMOS
% ---

% Resumo em português
% conter no máximo 500 palavras
\textoresumo{
    Este trabalho é um breve modelo  para a escrita de monografias de qualificação, dissertações e teses utilizando o ambiente \LaTeX, de acordo com as normas exigidas pelo Instituto de Ciências Matemáticas e de Computação (ICMC), da Universidade de São Paulo (USP). Para a confecção deste modelo foi utilizado a última versão (1.9.2) do pacote de classes \textit{abnTeX2} que segue as normas da Associação Brasileira de Normas Técnicas. A elaboração de uma monografia, dissertação ou tese pode ser feita sobrescrevendo o conteúdo deste modelo. 
    }{Modelo, Monografia de qualificação, Dissertação, Tese, Latex}

% ---
% resumo em inglês
% ---
\textoresumo[english]{
    This paper is a brief model for writing qualification monographs, dissertations and thesis using \LaTeX environment, in accordance with the standards required by the Institute of Mathematics and Computer Sciences (ICMC), University of São Paulo (USP). For making this model, the latest version (1.9.2) \textit{abnTeX2} classes package was used. This package follow the rules of the Brazilian Association of Technical Standards. A drafting a monograph, dissertation or thesis can be done by overwriting the contents of this model.
    }{Template, Qualification monograph, Dissertation, Thesis, Latex}
    
    
% ---
% Configurações de aparência do PDF final
% ---
% alterando o aspecto da cor azul
\definecolor{blue}{RGB}{41,5,195}

% informações do PDF
\makeatletter
\hypersetup{
     	pagebackref=true,
		pdftitle={\@title}, 
		pdfauthor={\@author},
    	pdfsubject={\imprimirpreambulo},
	    pdfcreator={LaTeX with abnTeX2/ICMC-USP},
		pdfkeywords={\palavraschave}, 
		colorlinks=true,       		% false: boxed links; true: colored links
    	linkcolor=blue,          	% color of internal links
    	citecolor=blue,        		% color of links to bibliography
    	filecolor=magenta,      	% color of file links
		urlcolor=blue,
		bookmarksdepth=4
}
\makeatother
% --- 

% ----------------------------------------------------------
% ELEMENTOS PRÉ-TEXTUAIS
% ----------------------------------------------------------

% Inserir a ficha catalográfica
% \incluifichacatalografica*{tex/fichaCatalografica.pdf}
 \incluifichacatalografica{634} % Código Cutter: número atribuído ao sobrenome do autor. Para obtê-lo, consulte a tabela Cutter Sanborn (em http://www.davignon.qc.ca/cutter1.html), procure pelo sobrenome ou forma mais próxima ao sobrenome completo e coloque o número indicado como parâmetro.


% DEDICATÓRIA / AGRADECIMENTO / EPÍGRAFE
\textodedicatoria*{tex/pre-textual/dedicatoria}
\textoagradecimentos*{tex/pre-textual/agradecimentos}
\textoepigrafe*{tex/pre-textual/epigrafe}

% Inclui a lista de figuras
\incluilistadefiguras

% Inclui a lista de tabelas
\incluilistadetabelas

% Inclui a lista de quadros
% \incluilistadequadros

% Inclui a lista de algoritmos
% \incluilistadealgoritmos

% Inclui a lista de códigos
% \incluilistadecodigos

% Inclui a lista de siglas e abreviaturas
\incluilistadesiglas

% Inclui a lista de símbolos
% \incluilistadesimbolos

% ----
% Início do documento
% ----
\begin{document}
% ----------------------------------------------------------
% ELEMENTOS TEXTUAIS
% ----------------------------------------------------------
\textual

\chapter{Introdução}
\label{chapter:introducao}
% Comando simples para exibir comandos Latex no texto
\newcommand{\comando}[1]{\textbf{$\backslash$#1}}

\section{Considerações Gerais}

Testando a inserção de uma nova sigla \sigla{EESC}{Escola de Engenharia de São Carlos}, vamos ver se deu certo.

\sigla{ESR}{Estrutura de Solo Reforçado}

\sigla{RCD}{Resíduo de Construção e Demolição}

\sigla{RCD-R}{Resíduo de Construção e Demolição Reciclado}


\begin{table}[htb]
\IBGEtab{%
  \caption{População dos países da América do Sul} \label{tabela:populacao_america_sul_nova}
}{%
\begin{tabular}{r|p{3cm}|r}        
\toprule
Código  & País            & População   \\ \midrule \midrule
1       & Brasil          & ~~~~~~191.480.630 \\ \midrule 
2       & Argentina       &  39.934.100 \\ \midrule 
3       & Colômbia        &  46.741.100 \\ \midrule 
4       & Paraguai        &   9.694.200 \\ \midrule 
5       & Uruguai         &   3.350.500 \\ \midrule 
6      & Chile           &  16.803.000 \\ \bottomrule
\end{tabular}
}{%
  \fdireta{wikipedia:2011:america_sul}%
  \nota{Esta é uma nota, que diz que os dados são baseados na
  regressão linear.}%
  \nota[Anotações]{Uma anotação adicional, que pode ser seguida de várias
  outras, porém são opcionais.}%
  }
\end{table}

\begin{table}[htb]
\caption{Editores de Texto Livres}
\label{quadro:editores_texto_livres_novo}
\centering
\begin{tabular}{|l|l|r|}        \hline
Editor     & Multiplataforma & Específico para Latex \\ \hline
Kwriter    & Sim             & Não                   \\
Texmaker   & Sim             & Sim                   \\
Kile       & Sim             & Sim                   \\
Geany      & Sim             & Não                   \\ \hline
\end{tabular}
\end{table}

\begin{figure}[htb]
 \caption{Logomarca da USP}
 \label{fig:logomarca_usp_novo}
 \centering
 \includegraphics[scale=0.3]{usp-logo.pdf}
 \fdireta{usp:logo}
\end{figure}



Este documento explica brevemente como trabalhar com a classe \LaTeX~\textit{icmc} para confeccionar trabalhos acadêmicos seguindo as normas da \sigla{ABNT}{Associação Brasileira de Normas Técnicas} e as \aspas{\textit{Diretrizes para apresentação de dissertações e teses da USP: documento eletrônico e impresso. Parte I (ABNT)}}, publicado pelo \sigla{SIBi}{Sistema Integrado de Bibliotecas} USP. O presente manual também atende as exigências prevista no regimento do Programa de Pós-graduação em \sigla{CCMC}{Ciências da Computação e Matemática Computacional} do \sigla{ICMC}{Instituto de Ciências Matemáticas e de Computação} da \sigla{USP}{Universidade de São Paulo}.


A classe \textit{icmc} foi construída com base na última versão da classe \textit{abntex2} e do pacote \textit{abntex2cite}. Portanto, este documento exemplifica a elaboração de trabalho
acadêmico (tese, dissertação e outros do gênero) produzido conforme a ABNT NBR
14724:2011 \textit{Informação e documentação - Trabalhos acadêmicos - Apresentação}.

Assim, é altamente recomendável que seja consultada a documentação do \textit{abntex2}\footnote{https://code.google.com/p/abntex2/downloads/list}. A classe \textit{abntex2} foi desenvolvida para facilitar a escrita de documentos seguindo as normas da ABNT no ambiente \LaTeX\;\cite{frasson:2005:classe_abnt}.

Todo o trabalho de pesquisa e ajustes da presente classe \LaTeX~\emph{icmc} foram feitos pelo aluno mestrado do Programa de Pós-graduação em Ciência da Computação e Matemática Computacional, Humberto Lidio Antonelli, durante a confecção da sua monografia de qualificação.

O requisito básico para utilização da classe \textit{icmc} é criar um documento desta classe com o comando
\comando{documentclass[@parameters]\{icmc\}} e ter, no diretório de trabalho, o arquivo \emph{icmc.cls} presente. Entretanto, recomenda-se fortemente manter a estrutura de diretório inicial fornecida por este modelo. Além disso, para que o documento esteja em conformidade com as normas exigidas pelo programa de Pós-Graduação, o \textbf{projeto deve ser compilado utilizando \textit{XeLaTeX} ou \textit{LuaLaTeX}}. Esse processo de compilação é necessário para que as fontes externas utilizadas para gerar a capa sejam incluídas.

Os parâmetros possíveis utilizados pelo \comando{documentclass} são:
\begin{description}
\item[qualificacao] Exclusivamente para monografias de qualificação em geral;
\item[mestrado / doutorado] Identifica o curso ao qual o aluno pertence, sendo utilizado apenas uma das duas opcões disponíveis. O valor padrão é \textbf{doutorado};
\item[pre-defesa / pos-defesa] Identifica a situação do documento (exceto para qualificação), sedo necessário apenas uma das duas opções. O valor padrão é \textbf{pos-defesa};
\item[french, spanish, english, brazil] Adiciona o idioma para correta hifenização correta no documento. O último idioma declarado é o principal do documento. O valor padrão é \textbf{brazil}.
\end{description}

\chapter{Instalando o abnTeX2}
\label{chapter:instalando-abntex}
A instalação do \emph{abnTeX2} varia de acordo com o sistema operacional empregado pelo usuário. Aqui serão apresentadas as formas de instalação nos sistemas mais utilizados atualmente, a saber: Linux (Ubuntu 12.04), Mac OS X e Windows 7

\section{Linux (Ubuntu 12.04)}

Se você já instalou o Tex Live via apt-get, basta seguir os seguintes comandos:

\begin{enumerate}

\item Baixe os arquivos de instalação do abnTeX2 (\url{http://code.google.com/p/abntex2/downloads/list}). Nesse link você também encontra a documentação e exemplos de uso.
\item Extraia o conteúdo do arquivo baixado na pasta texmf local, geralmente /usr/local/share/texmf. 
\item Em um Terminal: extraia o ZIP: \emph{unzip abntex2.tds.zip} em qualquer local;
\item copie o conteúdo extraído para o destino: \emph{cp abntex2/* /usr/local/share/texmf};
\item Em um Terminal digite: \emph{sudo texhash}
\item Pronto!
\end{enumerate}

\section{Mac OS}

Primeiramente, deve-se abrir o terminal do Mac que pode ser encontrado em Aplicativos/Utilitários - buscando pelo Finder.  E seguir os comandos abaixo:
\begin{enumerate}
\item Baixe os arquivos de instalação do abnTeX2 (\url{http://code.google.com/p/abntex2/downloads/list}). Nesse link você também encontra a documentação e exemplos de uso.
\item Extraia o conteúdo do arquivo baixado na pasta \emph{texmf} local, geralmente \emph{/usr/local/texlive/texmf-local}
\item Em um Terminal digite: \emph{sudo texhash}
\item Pronto!
\end{enumerate}
 
 \section{Windows 7}

\subsection{Instalar/atualizar pelo Package Manager (recomendado)}

Geralmente o abnTeX2 é baixado e instalado automaticamente pelo MiKTeX quando o usuário compila pela primeira vez um dos modelos do abnTeX2. Porém, caso isso não ocorra, siga os passos seguintes:

\begin{enumerate}
\item Clique em Iniciar/Start -> Todos os Programas/All Programs -> MiKTeX -> Package Manager;
\item Clique em Repository / Synchronize;
\item Clique com o botão direito sobre \emph{abntex2} na lista e selecione Install (ou Update, caso já esteja instalado);
\item Pronto!
\end{enumerate}

\subsection{Instalar/atualizar manualmente}

Você apenas precisará utilizar a instalação manual no caso de:

\begin{enumerate}
\item o abnTeX2 não estar na lista de pacotes do MiKTeX por alguma razão;
\item você não poder utilizar uma conexão com a Internet no momento da instalação;
\item a versão do abnTeX2 no MiKTeX estar desatualizada em relação à versão disponível no CTAN.
\end{enumerate}
Em qualquer caso, lembre-se de remover uma eventual instalação anterior do abnTeX2 . Se houver instalado pelo Package Manager, remova o abnTeX2 também por ele.

Passos para instalação manual do abnTeX2 no MiKTeX:

\begin{enumerate}
\item Baixe os arquivos de instalação do abnTeX2 (abntex2.tds-vX.X.zip). Nesse link você também encontra a documentação e exemplos de uso.
\item Extraia o conteúdo do arquivo baixado em uma pasta qualquer;
\item Você pode criar uma pasta abntex2, por exemplo, em $C:\backslash abntex2\backslash$;
\item Consulte http://www.tex.ac.uk/cgi-bin/texfaq2html?label=install-where para outras informações;
\item Clique em Iniciar/Start -> Todos os Programas/All Programs -> MiKTeX -> Settings;
\item Na aba Roots, adicione o diretório recém criado;
\item Na aba General, clique em Refresh FNDB, OU, se preferir, em um Terminal digite initexmf --update-fndb;
\item Pronto!
\end{enumerate}


% ---
% Finaliza a parte no bookmark do PDF, para que se inicie o bookmark na raiz
% ---
\bookmarksetup{startatroot}% 
% ---

% ----------------------------------------------------------
% ELEMENTOS PÓS-TEXTUAIS
% ----------------------------------------------------------
\postextual

% ----------------------------------------------------------
% Referências bibliográficas
% ----------------------------------------------------------
\bibliography{references}

% ----------------------------------------------------------
% Apêndices
% ----------------------------------------------------------

% ---
% Inicia os apêndices
% ---
\begin{apendicesenv}

    \chapter{Documento básico usando a classe \textit{icmc}}
    \label{chapter:documento-basico}
    \definecolor{gray}{rgb}{0.4,0.4,0.4}
\definecolor{darkblue}{rgb}{0.0,0.0,0.6}
\definecolor{cyan}{rgb}{0.0,0.6,0.6}
\definecolor{maroon}{rgb}{0.5,0,0}
\definecolor{darkgreen}{rgb}{0,0.5,0}


\lstdefinelanguage{myLatex}
{
    keywords={\titulo},
    alsoletter={-},
    sensitive=false,
    morecomment=[l]{\%},
    morecomment=[s]{/*}{*/},
    morestring=[b]",
    morestring=[b]',
    keywordstyle=\bfseries\color{blue},
    commentstyle=\itshape\color{darkgreen},
    morekeywords={documentclass, titulo, autor, data, orientador, coorientador, curso, textoresumo, incluifichacatalografica, textodedicatoria*, textoagradecimentos*, textoepigrafe*, incluilistadefiguras, incluilistadetabelas, incluilistadequadros, incluilistadealgoritmos, incluilistadecodigos, incluilistadesiglas, incluilistadesimbolos, textual, chapter, postextual, begin, bibliography, end}, 
alsoletter={*, \{, \}, \[, \]},
 morekeywords=[2]{\{, \}, \[, \]},
 keywordstyle=[2]\bfseries\color{blue},
 moredelim=[s][\color{maroon}]{\{}{\}},
    moredelim=[s][\itshape\color{maroon}]{\[}{\]},
}

%\lstdefinelanguage{TeX}
%{
%moredelim=*[s][\color{maroon}]{\{}{\}}
%otherkeywords={\{, \}, \[, \], \\}
%  morestring=[b]",
%  moredelim=[s][\bfseries\color{maroon}]{<}{\ },
%  moredelim=[s][\bfseries\color{maroon}]{</}{>},
%  moredelim=[l][\bfseries\color{maroon}]{/>},
%  moredelim=[l][\bfseries\color{maroon}]{>},
%  commentstyle=\color{darkgreen},
%  stringstyle=\color{blue},
%  identifierstyle=\color{red},
%  keywordstyle=\bfseries\color{maroon}
%moredelim=[l][\bfseries\color{maroon}]{>},
%commentstyle=\color{darkgreen},
%  stringstyle=\color{blue},
%  identifierstyle=\color{red}, moredelim=[l][\bfseries\color{maroon}]{\{},
%  keywordstyle=\bfseries\color{maroon}
%}

%\lstset{language={[LaTeX]TeX},
%texcsstyle=*\bfseries\color{blue},
%keywordstyle=\bfseries\color{blue},
%commentstyle=\color{darkgreen},
%morecomment=[s][\color{red}]{\{}{\}},
%otherkeywords={$, \{, \}, \[, \]}
%}

%\begin{codigo}[caption={Exemplo de um documento básico}, label={codigo:documento-basico}, language={[LaTeX]TeX},  breaklines=true,morekeywords={titulo, autor, data, orientador, coorientador, curso, textoresumo, incluifichacatalografica, textodedicatoria*, textoagradecimentos*, textoepigrafe*, incluilistadefiguras, incluilistadetabelas, incluilistadequadros, incluilistadealgoritmos, incluilistadecodigos, incluilistadesiglas, incluilistadesimbolos, {\backslash}textual, chapter, postextual}, alsoletter={{\backslash},*},morecomment=[s][\color{red}]{\{}{\}}]
\begin{codigo}[caption={Exemplo de um documento básico}, label={codigo:documento-basico}, language={myLatex},  breaklines=true]
% Documento utilizando a classe icmc
% Opções: 
%   Qualificação         = qualificacao 
%   Curso                = doutorado/mestrado
%   Situação do trabalho = pre-defesa/pos-defesa (exceto para qualificação)
% -- opções do pacote babel --
% Idioma padrão = brazil
	%spanish,			% idioma adicional para hifenização
	%english,			% idioma adicional para hifenização
	%brazil				% o último idioma é o principal do documento
\ documentclass[doutorado, spanish, english, brazil]{packages/icmc}

% Título do trabalho
\titulo{Título da Monografia}

% Nome do autor
\autor[Abreviação]{Nome completo do autor}

% Data do depósito
\data{18}{12}{2012}

% Nome do Orientador
\orientador[Orientador]{Titulação do orientador}{Nome completo do Orientador}

% Nome do Coorientador (caso não exista basta remover)
\coorientador[Coorientador]{Titulação do coorientador}{Nome completo do Coorientador}
% Se coorientadora troque Coorientador: por Coorientadora dentro do colchetes

% Sigla do programa de Pós-graduação (CCMC, MAT, PIPGES, PROFMAT, MECAI)
\curso{CCMC}
% O valor entre colchetes é opcional para este programa

% Resumo
\textoresumo[Idioma]{
Texto do resumo do trabalho.
}{Lista de palavras-chave separada por virgulas}

% ----------------------------------------------------------
% ELEMENTOS PRÉ-TEXTUAIS
% ----------------------------------------------------------

% Inserir a ficha catalográfica
\incluifichacatalografica{634} % Código Cutter: número atribuído ao sobrenome do autor (disponível em http://www.davignon.qc.ca/cutter1.html).

% Incluí o texto da Dedicatória
\textodedicatoria*{tex/pre-textual/dedicatoria}

% Incluí o texto dos Agradecimentos
\textoagradecimentos*{tex/pre-textual/agradecimentos}

% Incluí o texto da Epígrafe
\textoepigrafe*{tex/pre-textual/epigrafe}

% Inclui a lista de figuras
\incluilistadefiguras

% Inclui a lista de tabelas
\incluilistadetabelas

% Inclui a lista de quadros
\incluilistadequadros

% Inclui a lista de algoritmos
\incluilistadealgoritmos

% Inclui a lista de códigos
\incluilistadecodigos

% Inclui a lista de siglas e abreviaturas
\incluilistadesiglas

% Inclui a lista de símbolos
\incluilistadesimbolos

% Início do documento
\begin{document}

% ----------------------------------------------------------
% ELEMENTOS TEXTUAIS
% ----------------------------------------------------------
\textual

\chapter{Introdução}

Capítulo de Introdução

\chapter{Desenvolvimento}

Capítulo de Desenvolvimento

\chapter{Conclusão}

Capítulo de conclusão

% ----------------------------------------------------------
% ELEMENTOS PÓS-TEXTUAIS
% ----------------------------------------------------------
\postextual

% Nome do arquivo com as referências bibliográficas
\bibliography{referencias}

\end{document}

\end{codigo}

\end{apendicesenv}
% ---


% ----------------------------------------------------------
% Anexos
% ----------------------------------------------------------

% ---
% Inicia os anexos
% ---
\begin{anexosenv}

    \chapter{Páginas interessantes na Internet} 
    \label{chapter:paginas-interessantes}
    \begin{description}
 \item[\url{http://www.tex-br.org}] Página em português com diversos tutoriais e referências interessantes sobre \LaTeX;
 \item[\url{http://en.wikibooks.org/wiki/LaTeX}] Livro em formato \textit{wiki} gratuito sobre \LaTeX;
 \item[\url{http://tobi.oetiker.ch/lshort/lshort.pdf}] Ótimo tutorial sobre \LaTeX (possui versão em português \url{http://alfarrabio.di.uminho.pt/~albie/lshort/ptlshort.pdf}, mas a versão em inglês é a mais atual);
 \item[\url{http://code.google.com/p/abntex2/}] Página do abnTeX2, grupo que desenvolve os pacotes e classes em \LaTeX para as normas da ABNT, nos quais a classe \textit{icmc} foi baseada;
\item[\url{http://www.more.ufsc.br}] Página do Mecanismo On-line para Referências  (MORE) desenvolvido pela UFSC;
\item[\url{http://detexify.kirelabs.org/classify.html}] Página para recuperar o código de símbolos em \LaTeX a partir do desenho fornecido pelo usuário.
 \end{description}

\end{anexosenv}
% ---

\end{document}